\documentclass[10pt]{article}
\usepackage{resume}
\hypersetup{pdftitle={Zhexiang Li - Research Resume},pdflang={en-US}}
\begin{document}
\englishheader

\section{Education}
\entry{University of California, Los Angeles}{Expected Jun 2028}
M.S. in Computer Science
\entry{University of Southern California}{May 2026}
B.S. in Computer Science; GPA: 3.76/4.00; Viterbi Dean's List

\section{Research Skills}
\textbf{Methods:} multimodal learning, efficient attention, token routing, molecular property prediction\\
\textbf{Tools:} Python, PyTorch, Triton, RDKit, Chemprop, C/C++, SQL, Linux

\section{Research Experience}
\entry{Graduate Researcher | Health Intelligence Lab, UCLA}{Aug 2026 - Present}
\detail{Advisor: Prof. Yuzhe Yang}
\begin{itemize}
\item Investigate multimodal glucose forecasting using CGM, insulin, meals, activity, and patient context.
\item Built patient-disjoint MetaboNet benchmarks and data audits to compare GluFormer, T1DFormer, and custom models at 30-, 60-, and 120-minute horizons.
\end{itemize}
\entry{Drug Discovery Research | USC, Prof. Yan Liu's group}{Spring 2026}
\detail{Research collaboration with a Ph.D. student in the group}
\begin{itemize}
\item Benchmarked molecular descriptor-based models, Chemprop, and hybrid approaches across seven Therapeutics Data Commons datasets using scaffold splits.
\item Analyzed descriptor importance and SHAP attributions to examine molecular features associated with predicted antiviral activity.
\end{itemize}

\section{Selected Papers \& Research Contributions}
\entry{MedTokenBudget: Lesion-Preserving Token Routing}{2026}
\detail{Zhexiang Li (sole author). Preprint. \link{https://arxiv.org/abs/2609.30613}{Paper}}
\begin{itemize}
\item Developed lesion-supervised post-backbone token routing to retain diagnostic image regions under a limited token budget.
\item At a 30\% token budget on ISIC 2019, achieved \textbf{80.38 $\pm$ 0.18\%} accuracy versus 76.93 $\pm$ 0.03\% for ToMe across three seeds; the frozen backbone remains unchanged.
\end{itemize}
\entry{Mass-Conservation Token Transformers for MS/MS Retrieval}{2026}
\detail{Manuscript submitted to AAAI 2027}
\begin{itemize}
\item Encoded mass-conserving fragment pairs as sparse evidence tokens for spectrum retrieval; constructed mass- and overlap-based hard negatives and evaluated ranking changes after evidence removal.
\end{itemize}
\entry{Efficient Attention via Pre-Scoring}{2025}
\detail{Zhexiang Li, Haoyu Wang, Yutong Bao, David Woodruff. Under review. \link{https://arxiv.org/abs/2505.11040}{Paper}}
\begin{itemize}
\item Proposed pre-scoring to prioritize informative keys before approximate attention; implemented key selection in Triton as first author, advised by Prof. David P. Woodruff (CMU).
\item Integrated with HyperAttention, reduced ChatGLM perplexity from \textbf{12.0 to 9.5} at 131k-token contexts under a fixed interaction budget.
\end{itemize}

\section{Teaching}
\entry{Machine Learning Tutor | UCSC}{Spring 2023}
\begin{itemize}
\item Led office hours and co-designed Jupyter assignments on logistic regression, SVMs, and tree models; contributed to the coding autograder for Prof. Lise Getoor's course.
\end{itemize}
\end{document}
